% GENERATED FILE. Edit canonical lessons, then run npm run generate:books.
% canonical-source-hash: fnv1a64:f7c8c30969dab69b
% canonical-lessons: KA-C350-sanmanisu, KA-C350-kondadu, KA-C350-birudu, KA-C350-kirti, KA-C350-udghatisu

\chapter{An Evening of Honours}
\label{ch:ka-an-evening-of-honours}
\hlchaptermodality{350}

\begin{chapteropening}
\textbf{By the end of this chapter:} \emph{I can say to honour, to extol, an honorific title, fame and to inaugurate.}

Complete the last lesson of chapter 350: I can say to honour, to extol, an honorific title, fame and to inaugurate.
\end{chapteropening}

\section[\texorpdfstring{\kn{ಸನ್ಮಾನಿಸು} (sanmānisu)}{ಸನ್ಮಾನಿಸು (sanmānisu)}]{\kn{ಸನ್ಮಾನಿಸು} (sanmānisu) — to honour; to felicitate}

\label{lesson:KA-C350-sanmanisu}



\begin{quote}
Before the new one: say the Kannada for a charioteer, then the Kannada for to fall in battle.
\end{quote}

\subsection*{You'll want to know: \kn{ಸನ್ಮಾನಿಸು}}
\textbf{\kn{ಸನ್ಮಾನಿಸು}} — \emph{sanmānisu} — \textquotedblleft{}to honour; to felicitate\textquotedblright{}.

\begin{quote}
\kn{ಊರಿನವರು} \kn{ಹಳೆಯ} \kn{ಶಿಕ್ಷಕರನ್ನು} \kn{ಸನ್ಮಾನಿಸಿದರು}.
\end{quote}

\begin{quote}
\emph{ūrinavaru haḷeya śikṣakarannu sanmānisidaru.}
\end{quote}

\begin{quote}
The townspeople honoured their old teacher.
\end{quote}

Past (I …): \kn{ಸನ್ಮಾನಿಸಿದೆ}, \emph{sanmāniside}.

The person honoured takes -\kn{ಅನ್ನು}, often with a shawl and a garland. Towns also honour the workers who keep their streets \kn{ಸ್ವಚ್ಛ}, clean.

\begin{cousinweb}
\textbf{\kn{ಸನ್ಮಾನಿಸು}} (\emph{sanmānisu}) is built on Sanskrit \emph{sanmāna}, \textquotedblleft{}respect or esteem\textquotedblright{} (Monier-Williams), with the Kannada verb ending -\emph{isu}.
\end{cousinweb}

\begin{grammarlens}[title={What you've built}]
One more word for telling what happened.
\end{grammarlens}

\subsection*{Your turn}
\begin{itemize}
\raggedright
  \item \emph{Say it:} \emph{sanmānisu}
  \item \emph{Say it:} \emph{sanmānisu}, once more
  \item \emph{Say it:} \emph{ūrinavaru haḷeya śikṣakarannu sanmānisidaru.} — The townspeople honoured their old teacher.
  \item \emph{Recall:} say the Kannada for a charioteer, then the Kannada for to fall in battle, then say \emph{sanmānisu} again
\end{itemize}

\begin{tcolorbox}[breakable,colback=teal!4,colframe=teal!35!black,title={Before you move on}]
What does \kn{ಸನ್ಮಾನಿಸು} mean? (\textquotedblleft{}To honour; to felicitate\textquotedblright{}.) Say it once more.
\end{tcolorbox}
\section[\texorpdfstring{\kn{ಕೊಂಡಾಡು} (koṇḍāḍu)}{ಕೊಂಡಾಡು (koṇḍāḍu)}]{\kn{ಕೊಂಡಾಡು} (koṇḍāḍu) — to extol; to praise highly}

\label{lesson:KA-C350-kondadu}



\begin{quote}
Before the new one: say the Kannada for to fall in battle, then the Kannada for to honour.
\end{quote}

\subsection*{You'll want to know: \kn{ಕೊಂಡಾಡು}}
\textbf{\kn{ಕೊಂಡಾಡು}} — \emph{koṇḍāḍu} — \textquotedblleft{}to extol; to praise highly\textquotedblright{}.

\begin{quote}
\kn{ಜನರು} \kn{ಆ} \kn{ಶಿಕ್ಷಕರ} \kn{ಸೇವೆಯನ್ನು} \kn{ಕೊಂಡಾಡಿದರು}.
\end{quote}

\begin{quote}
\emph{janaru ā śikṣakara sēveyannu koṇḍāḍidaru.}
\end{quote}

\begin{quote}
People praised that teacher's service.
\end{quote}

Past (I …): \kn{ಕೊಂಡಾಡಿದೆ}, \emph{koṇḍāḍide}.

Stronger and more public than plain praise: to sing someone's praises.

\begin{cousinweb}
\textbf{\kn{ಕೊಂಡಾಡು}} (\emph{koṇḍāḍu}) is usually explained as a native compound of \kn{ಕೊಂಡು}, from \emph{koḷ}, \textquotedblleft{}to take, to accept\textquotedblright{} (DEDR 2151), and \emph{āḍu}, \textquotedblleft{}to act, to speak\textquotedblright{} (DEDR 347).
\end{cousinweb}

\begin{grammarlens}[title={What you've built}]
One more word for telling what happened.
\end{grammarlens}

\subsection*{Your turn}
\begin{itemize}
\raggedright
  \item \emph{Say it:} \emph{koṇḍāḍu}
  \item \emph{Say it:} \emph{koṇḍāḍu}, once more
  \item \emph{Say it:} \emph{janaru ā śikṣakara sēveyannu koṇḍāḍidaru.} — People praised that teacher's service.
  \item \emph{Recall:} say the Kannada for to fall in battle, then the Kannada for to honour, then say \emph{koṇḍāḍu} again
\end{itemize}

\begin{tcolorbox}[breakable,colback=teal!4,colframe=teal!35!black,title={Before you move on}]
What does \kn{ಕೊಂಡಾಡು} mean? (\textquotedblleft{}To extol; to praise highly\textquotedblright{}.) Say it once more.
\end{tcolorbox}
\section[\texorpdfstring{\kn{ಬಿರುದು} (birudu)}{ಬಿರುದು (birudu)}]{\kn{ಬಿರುದು} (birudu) — an honorific title}

\label{lesson:KA-C350-birudu}



\begin{quote}
Before the new one: say the Kannada for to honour, then the Kannada for to extol.
\end{quote}

\subsection*{You'll want to know: \kn{ಬಿರುದು}}
\textbf{\kn{ಬಿರುದು}} — \emph{birudu} — \textquotedblleft{}an honorific title\textquotedblright{}.

\begin{quote}
\kn{ರಾಜ} \kn{ಆ} \kn{ಕವಿಗೆ} \kn{ಒಂದು} \kn{ಬಿರುದು} \kn{ಕೊಟ್ಟನು}.
\end{quote}

\begin{quote}
\emph{rāja ā kavige ondu birudu koṭṭanu.}
\end{quote}

\begin{quote}
The king gave that poet an honorific title.
\end{quote}

Plural: \kn{ಬಿರುದುಗಳು}, \emph{birudugaḷu}.

A title of honour given by a king or a society, used with the name.

\begin{cousinweb}
\textbf{\kn{ಬಿರುದು}} (\emph{birudu}) is native Dravidian: DEDR 5414 lists Kannada \emph{biridu}, \emph{biruda}, \emph{birudu}, \textquotedblleft{}praise, a badge of honour\textquotedblright{}, beside Tamil \emph{virutu}, \textquotedblleft{}a title, a banner, a badge of victory\textquotedblright{}; Sanskrit has a similar \emph{biruda}.
\end{cousinweb}

\begin{grammarlens}[title={What you've built}]
One more word for telling what happened.
\end{grammarlens}

\subsection*{Your turn}
\begin{itemize}
\raggedright
  \item \emph{Say it:} \emph{birudu}
  \item \emph{Say it:} \emph{birudu}, once more
  \item \emph{Say it:} \emph{rāja ā kavige ondu birudu koṭṭanu.} — The king gave that poet an honorific title.
  \item \emph{Recall:} say the Kannada for to honour, then the Kannada for to extol, then say \emph{birudu} again
\end{itemize}

\begin{tcolorbox}[breakable,colback=teal!4,colframe=teal!35!black,title={Before you move on}]
What does \kn{ಬಿರುದು} mean? (\textquotedblleft{}An honorific title\textquotedblright{}.) Say it once more.
\end{tcolorbox}
\section[\texorpdfstring{\kn{ಕೀರ್ತಿ} (kīrti)}{ಕೀರ್ತಿ (kīrti)}]{\kn{ಕೀರ್ತಿ} (kīrti) — fame; glory}

\label{lesson:KA-C350-kirti}



\begin{quote}
Before the new one: say the Kannada for to extol, then the Kannada for an honorific title.
\end{quote}

\subsection*{You'll want to know: \kn{ಕೀರ್ತಿ}}
\textbf{\kn{ಕೀರ್ತಿ}} — \emph{kīrti} — \textquotedblleft{}fame; glory\textquotedblright{}.

\begin{quote}
\kn{ಆ} \kn{ಆಟಗಾರನ} \kn{ಕೀರ್ತಿ} \kn{ದೇಶದ} \kn{ತುಂಬಾ} \kn{ಹರಡಿತು}.
\end{quote}

\begin{quote}
\emph{ā āṭagārana kīrti dēśada tumbā haraḍitu.}
\end{quote}

\begin{quote}
That player's fame spread across the country.
\end{quote}

\kn{ಕೀರ್ತಿ} \kn{ತರು}, to bring glory: \kn{ಅವಳು} \kn{ಊರಿಗೆ} \kn{ಕೀರ್ತಿ} \kn{ತಂದಳು}.

\begin{cousinweb}
\textbf{\kn{ಕೀರ್ತಿ}} (\emph{kīrti}) is borrowed from Sanskrit \emph{kīrti}, \textquotedblleft{}report\textquotedblright{}, and later \textquotedblleft{}fame\textquotedblright{} (CDIAL 3200).
\end{cousinweb}

\begin{grammarlens}[title={What you've built}]
One more word for telling what happened.
\end{grammarlens}

\subsection*{Your turn}
\begin{itemize}
\raggedright
  \item \emph{Say it:} \emph{kīrti}
  \item \emph{Say it:} \emph{kīrti}, once more
  \item \emph{Say it:} \emph{ā āṭagārana kīrti dēśada tumbā haraḍitu.} — That player's fame spread across the country.
  \item \emph{Recall:} say the Kannada for to extol, then the Kannada for an honorific title, then say \emph{kīrti} again
\end{itemize}

\begin{tcolorbox}[breakable,colback=teal!4,colframe=teal!35!black,title={Before you move on}]
What does \kn{ಕೀರ್ತಿ} mean? (\textquotedblleft{}Fame; glory\textquotedblright{}.) Say it once more.
\end{tcolorbox}
\section[\texorpdfstring{\kn{ಉದ್ಘಾಟಿಸು} (udghāṭisu)}{ಉದ್ಘಾಟಿಸು (udghāṭisu)}]{\kn{ಉದ್ಘಾಟಿಸು} (udghāṭisu) — to inaugurate; to open}

\label{lesson:KA-C350-udghatisu}



\begin{quote}
Before the new one: say the Kannada for an honorific title, then the Kannada for fame.
\end{quote}

\subsection*{You'll want to know: \kn{ಉದ್ಘಾಟಿಸು}}
\textbf{\kn{ಉದ್ಘಾಟಿಸು}} — \emph{udghāṭisu} — \textquotedblleft{}to inaugurate; to open\textquotedblright{}.

\begin{quote}
\kn{ಹಿರಿಯ} \kn{ಲೇಖಕರು} \kn{ಹೊಸ} \kn{ಗ್ರಂಥಾಲಯವನ್ನು} \kn{ಉದ್ಘಾಟಿಸಿದರು}.
\end{quote}

\begin{quote}
\emph{hiriya lēkhakaru hosa granthālayavannu udghāṭisidaru.}
\end{quote}

\begin{quote}
A senior writer opened the new library.
\end{quote}

Past (I …): \kn{ಉದ್ಘಾಟಿಸಿದೆ}, \emph{udghāṭiside}.

Buildings, events and meetings are opened: \kn{ಸಮಾರಂಭವನ್ನು} \kn{ಉದ್ಘಾಟಿಸು}. The noun is \kn{ಉದ್ಘಾಟನೆ}.

\begin{cousinweb}
\textbf{\kn{ಉದ್ಘಾಟಿಸು}} (\emph{udghāṭisu}) is built on Sanskrit \emph{udghāṭana}, \textquotedblleft{}the act of opening, unlocking\textquotedblright{} (Monier-Williams), with the Kannada verb ending -\emph{isu}.
\end{cousinweb}

\begin{grammarlens}[title={What you've built}]
One more word for telling what happened.
\end{grammarlens}

\subsection*{Your turn}
\begin{itemize}
\raggedright
  \item \emph{Say it:} \emph{udghāṭisu}
  \item \emph{Say it:} \emph{udghāṭisu}, once more
  \item \emph{Say it:} \emph{hiriya lēkhakaru hosa granthālayavannu udghāṭisidaru.} — A senior writer opened the new library.
  \item \emph{Recall:} say the Kannada for an honorific title, then the Kannada for fame, then say \emph{udghāṭisu} again
\end{itemize}

\begin{tcolorbox}[breakable,colback=teal!4,colframe=teal!35!black,title={Before you move on}]
What does \kn{ಉದ್ಘಾಟಿಸು} mean? (\textquotedblleft{}To inaugurate; to open\textquotedblright{}.) Say it once more.
\end{tcolorbox}
